% Generated draft: requires mathematical and source review.
\documentclass[11pt]{article}
\usepackage[margin=25mm]{geometry}
\usepackage{fontspec}
\setmainfont{Latin Modern Roman}
\usepackage{amsmath,amssymb,mathtools,unicode-math}
\setmathfont{Latin Modern Math}
\usepackage{enumitem,needspace,xurl,hyperref}
\setlength{\parindent}{0pt}
\setlength{\parskip}{5pt}
\setlength{\emergencystretch}{3em}
\newcommand{\N}{\mathbb{N}}
\newcommand{\Z}{\mathbb{Z}}
\newcommand{\Q}{\mathbb{Q}}
\newcommand{\R}{\mathbb{R}}
\newcommand{\C}{\mathbb{C}}
\newcommand{\F}{\mathbb{F}}
\newcommand{\A}{\mathbb{A}}
\newcommand{\PP}{\mathbb{P}}
\newcommand{\E}{\mathbb{E}}
\newcommand{\eps}{\varepsilon}
\newcommand{\dd}{\,\mathrm{d}}
\begin{document}
\section*{UnsolvedMath ID: 30002663}
\textbf{Source problem code:} catalogue entry
\subsection{Definitions and mathematical statement}
Equip \(\R^{2n}\), for a positive integer \(n\), with the standard symplectic form \(\omega_0=\sum_{j=1}^{n}\dd x_j\wedge\dd y_j\). A convex body means a compact convex subset with nonempty interior. A finite convex cover of a convex body \(X\) consists of convex bodies \(X_1,\ldots,X_m\subset\R^{2n}\), where \(m\geq1\), satisfying
\[
X\subseteq\bigcup_{i=1}^{m}X_i.
\]
The covering bodies are not required to lie inside \(X\).

In the usual normalized terminology, a symplectic capacity is a function \(c\) on a specified class of symplectic sets, with values in \([0,\infty]\), which is monotone under symplectic embeddings, satisfies \(c(\lambda A)=\lambda^2c(A)\) for \(\lambda>0\), and obeys
\[
c(B^{2n}(1))=c(Z^{2n}(1))=\pi,
\]
where
\[
B^{2n}(1)=\left\{\sum_{j=1}^{n}(x_j^2+y_j^2)<1\right\},
\qquad
Z^{2n}(1)=\{x_1^2+y_1^2<1\}.
\]
Here a symplectic embedding \(\phi\) satisfies \(\phi^*\omega_0=\omega_0\). These are standard conventions for explaining the term; the supplied record does not specify the capacity's domain, normalization, or how its values on closed convex bodies are defined from any domain of open sets. In particular, using \(c(X)\) directly or interpreting it as \(c(\operatorname{int}X)\) requires a convention not supplied in the record.

\textbf{Subadditivity question, with the literal quantifier order of the supplied statement.} For every positive integer \(n\), every convex body \(X\subset\R^{2n}\), and every finite convex cover \(X_1,\ldots,X_m\) of \(X\), does there exist a symplectic capacity \(c\), defined on the sets involved, such that
\[
c(X)\leq\sum_{i=1}^{m}c(X_i)?
\]
With this quantifier order, \(c\) may depend on the cover. The record does not clarify whether the intended conjecture instead asks for one capacity satisfying this inequality for every finite convex cover in a fixed dimension, or for a compatible capacity across dimensions. Those are stronger uniformity requirements and are not silently imposed here. The unspecified capacity category and boundary convention remain qualifications requiring review.
\subsection{Short English statement}
The supplied question asks whether each finite convex cover of a convex body admits a symplectic capacity for which the sum of the capacities of the covering bodies is at least the capacity of the covered body. It remains unclear whether one capacity must work for all covers and how capacities are evaluated on closed bodies. The record labels the question open and describes it as a conjectural ingredient for Bang's problem; these are supplied assessments.
\subsection{Sources}
\begin{itemize}
\item Supplied local UnsolvedMath registry record, permanent ID \(30002663\), source problem code \texttt{catalogue entry}, title: \emph{Subadditivity of Symplectic Capacities Under Convex Covers}. Collection: Oberwolfach Reports Open Problems. No individual proposer is supplied; the proposed year is recorded as 2014.
\item Supplied report citation: \emph{Discrete Geometry (2015)}, Oberwolfach Reports, DOI 10.4171/owr/2014/40, pp. 2235--2296. Original-report link: \url{https://doi.org/10.4171/owr/2014/40}. The citation's 2015 label and the record's 2014 proposed year and DOI year are retained without resolving their relationship.
\item Supplied literature reference: Akopyan, Karasev, and Petrov, \emph{Bang problem}, \url{https://arxiv.org/abs/1404.0871}.
\item Literature assessment checked at \texttt{2026-08-21}: supplied dated curation, not independently verified. It labels the problem open, describes the subadditivity principle as a conjectural symplectic ingredient for Bang's problem, and reports that no general proof or disproof was found. The background explicitly requests independent verification of time-sensitive status and cited claims.
\item The registry labels the statement \texttt{corrected\_verified} and says that notation and publication prose were repaired without changing its substance. This is supplied editorial metadata, not independent mathematical or source verification of this draft.
\end{itemize}
\end{document}
